% =====================================================================
%  Workshop paper draft -- Janus PEFT modularity
%  Self-contained article (compiles on Overleaf with references.bib + figs/).
%  For submission: swap the preamble for neurips_2025.sty; body unchanged.
%  NUMBERS are verbatim from our experiments / the literature. "[ongoing]"
%  marks results not yet available. We report only consolidated findings.
% =====================================================================
\documentclass{article}
\usepackage[margin=2.3cm]{geometry}
\usepackage{hyperref}
\usepackage{booktabs}
\usepackage{amsmath}
\usepackage{amssymb}
\usepackage{natbib}
\usepackage{graphicx}
\usepackage{subcaption}
\usepackage{float}
\usepackage{titlesec}

% --- compact vertical spacing ---
\setlength{\parskip}{0.3ex}
\titlespacing*{\section}{0pt}{1.3ex}{0.6ex}
\titlespacing*{\subsection}{0pt}{0.9ex}{0.4ex}
\titlespacing*{\paragraph}{0pt}{0.5ex}{0.6em}
\titleformat{\section}{\normalfont\large\bfseries}{\thesection}{0.6em}{}
\titleformat{\subsection}{\normalfont\normalsize\bfseries}{\thesubsection}{0.5em}{}

\title{\vspace{-2.5em}\textbf{Parameter-Efficient Fine-Tuning Modularity in Unified Multimodal Models}\vspace{-1em}}
\author{}
\date{}

\begin{document}
\maketitle

\begin{abstract}
Unified multimodal models perform both image understanding and image generation
within a single shared transformer. This raises a question that does not exist
for a task specialist: can a parameter-efficient adapter specialize one
capability without degrading the others? We use style adaptation as a probe and
study the modularity of LoRA fine-tuning on Janus-Pro-7B, a unified model with
decoupled understanding and generation encoders over a shared backbone.
Fine-tuning the image-generation direction on a distant visual style leaves
language modelling (WikiText perplexity), factual knowledge (MMLU) and visual
understanding (POPE, MME) unchanged, across LoRA ranks 2 to 64 and across
KL-to-base regularization strengths. Fine-tuning the text direction on three
targets of increasing departure from the base -- a poetic style, a medical
domain, and a different language (Welsh) -- acquires each while leaving knowledge
and understanding unchanged; Welsh, the hardest, is acquired most strongly
(held-out perplexity $-64\%$) yet with the smallest weight change of the three.
These are consistent signals that PEFT on a unified model can specialize one
direction at little measurable cross-capability cost on standard benchmarks. We
outline the generative-quality measurements still needed to further delimit where
this modularity holds.
\end{abstract}

\section{Introduction}
Unified multimodal models process text and images in a single shared sequence
and serve both understanding (image$+$text to text) and generation (text to
image) from one backbone \citep{chameleon2024,wu2024janus,lu2024unifiedio2}.
Parameter-efficient fine-tuning, and LoRA in particular \citep{hu2022lora}, is
the standard way to specialize such a model on a downstream target at a fraction
of the cost of full fine-tuning.

This pairing raises a question that is meaningful only for a unified model. When
PEFT specializes one capability direction, what happens to the others? For a task
specialist, such as a text-to-image diffusion model, the question is undefined,
because there are no other modalities to preserve. For a unified model it is the
natural test of the ``one backbone, many capabilities'' premise. We therefore
treat style adaptation not as a goal (a specialist would transfer style more
cheaply and sharply) but as a probe with which to characterize the modularity of
PEFT: whether, and how locally, specializing one direction leaves the others
intact.

Prior work bounds the problem without answering it directly. LoRA forgets less
than full fine-tuning while also learning less \citep{biderman2024lora}, and
adapter capacity past a small rank yields diminishing returns
\citep{zhang2024scalingmeetsfinetuning}. In unified models, methods that
reconcile the understanding/generation tension typically modify the architecture
\citep{dou2024loramoe}. What is missing is a direct characterization of plain
LoRA's cross-capability footprint, measured in all directions, with a fine-tuning
target whose distance from the base is controlled so that ``no forgetting'' is
not an artifact of a target the model already mastered.

\paragraph{Contributions.}
(i) We pose a question specific to unified models, whether a parameter-efficient
adapter can specialize one capability direction without degrading the others, and
give a protocol to answer it: each fine-tuning target is chosen by its distance
from the base (poetry, $4.0\times$ the perplexity of a prose control; a narrow
anime style the base renders weakly), so ``no forgetting'' cannot be an artifact of
a target the model already mastered, and after specializing one direction we probe
\emph{all} the others (language modelling, knowledge, visual understanding).
(ii) Across both the generation and text directions -- the latter probed on three
targets of increasing departure from the base (a poetic style, a medical domain,
and a new language, Welsh) -- and robustly across LoRA rank (up to $64$) and
KL-to-base strength ($\lambda\in[0,1.5]$), specialization leaves the non-target
directions at base level on public benchmarks (MMLU, POPE, MME, WikiText), while
the target is acquired, measured for images with robust distributional metrics
(CMMD, KID, DINOv2, CLIP-IQA, directional similarity) rather than log-likelihood
alone.
(iii) We rule out the objection that this modularity is trivial because the adapter
barely changes the model: the induced weight update
$\lVert\Delta W\rVert_F/\lVert W\rVert_F$ is small ($0.6$--$1.9\%$) but real and
ordered (growing with rank, shrinking with the KL weight in both directions) and
\emph{modality-localized}, concentrated in the MLP with attention $q/k/v$ largely
spared, on different sub-modules for the image and text directions. A targeted,
low-norm, localized perturbation suffices.

\section{Related Work}
\paragraph{PEFT and forgetting.}
Adapters \citep{houlsby2019parameter} and LoRA \citep{hu2022lora} update a small
fraction of parameters. LoRA's adaptation/forgetting trade-off is comparatively
benign: it learns less and forgets less than full fine-tuning
\citep{biderman2024lora}, and rank past a small value gives diminishing returns
\citep{zhang2024scalingmeetsfinetuning}. A complementary line mitigates forgetting
by routing among multiple LoRAs \citep{dou2024loramoe,wu2024mole,li2024mixlora}.
These works study text-only or late-fusion models and treat forgetting within a
single modality.

\paragraph{Native and unified multimodal models.}
Fully-tokenized native models such as Chameleon \citep{chameleon2024}, the 4M
family \citep{bachmann20244m21} and Unified-IO\,2 \citep{lu2024unifiedio2} process
all modalities in one sequence. Janus and Janus-Pro \citep{wu2024janus,januspro}
decouple the visual encoders for understanding and generation while keeping a
single shared backbone. PEFT has been applied to such models
\citep{chern2024anole,serra2024narrowgate,wang2025vora}, but as a means to an end
rather than to characterize cross-capability interference. We neither modify the
architecture nor propose a new adapter; we measure the modularity of plain LoRA
across the directions of an existing unified model.

\section{Experimental Setup}
\paragraph{Model, directions, and adapter.}
We use Janus-Pro-7B \citep{januspro}: a SigLIP encoder for understanding and a
VQ-VAE (codebook $16{,}384$, $576$ tokens per image) for generation, over one
shared LLaMA-style transformer. It exposes three capability directions:
text-to-text language modelling, image$+$text understanding, and text-to-image
generation. We inject LoRA into the shared backbone
(\texttt{q,k,v,o,gate,up,down}); the generation forward traverses only the
language-model tensors, while understanding additionally engages SigLIP.

\paragraph{KL-to-base regularizer.}
For both directions we optionally anchor the adapter to the frozen base. Writing
$p_{\text{adapter}}(\cdot\mid x_{<t})$ for the adapted next-token distribution and
$p_{\text{base}}(\cdot\mid x_{<t})$ for the \emph{same} model with the LoRA adapter
disabled, the training loss is
\begin{equation}
\mathcal{L}=
\underbrace{\mathrm{CE}\big(p_{\text{adapter}},\,y\big)}_{\text{acquire the target}}
\;+\;\lambda\,
\underbrace{\frac{1}{T}\sum_{t=1}^{T}
\mathrm{KL}\!\Big(p_{\text{adapter}}(\cdot\mid x_{<t})\,\big\|\,
\mathrm{sg}\big[p_{\text{base}}(\cdot\mid x_{<t})\big]\Big)}_{\text{stay close to the base}}.
\label{eq:kl}
\end{equation}
The cross-entropy pulls the LoRA adapter toward the fine-tuning target; the KL
term, averaged over the $T$ token positions, penalizes drift of the adapter's
output distribution from the base's on the same inputs. The reference
$p_{\text{base}}$ is obtained for free by disabling the adapter, no separate
teacher model, and is held constant through a stop-gradient $\mathrm{sg}[\cdot]$,
so the KL moves the adapter toward the base and not the reverse. The single weight
$\lambda$ trades acquisition for preservation: $\lambda{=}0$ is plain
fine-tuning, larger $\lambda$ keeps the adapter near the base. It is the knob
swept in Tables~\ref{tab:kl} and~\ref{tab:text-kl}.

\paragraph{Why this regularizer.}
Anchoring a fine-tuned model to a frozen reference is the standard defence against
forgetting, realised either in weight space---L2-SP \citep{xuhong2018l2sp}, or its
Fisher-weighted form EWC \citep{kirkpatrick2017ewc}---or in function space---%
knowledge distillation \citep{hinton2015distilling}, Learning-without-Forgetting
\citep{li2017lwf}, and the KL penalty that keeps an RLHF policy near its reference
\citep{stiennon2020,ouyang2022instructgpt}. We adopt the function-space KL of
Eq.~\eqref{eq:kl} because the quantity we wish to preserve is the model's output
\emph{distribution}, for which the KL to the base is the natural divergence: a
weight-space penalty constrains parameters regardless of their functional effect,
while a logit MSE is sensitive to the (uncalibrated) logit scale. An ablation against L2-SP, logit-MSE and the reconstruction loss of Anole \citep{chern2024anole} at rank 32 (Appendix~\ref{app:reg-ablation}) supports this: under a trigger-free evaluation and at matched intervention strength, the KL attains the best style acquisition of all methods (CMMD $.108$ vs.\ $.155$ for no regularizer), while logit-MSE over-constrains the output and fails to transfer the style, and Anole acquires it but without anchoring to the base.



\paragraph{Forgetting and target selection.}
We probe the non-target directions with four public benchmarks scored by
log-likelihood only (no free-form decoding, hence deterministic and judge-free):
WikiText-2 perplexity \citep{wikitext}, MMLU \citep{mmlu}, POPE \citep{pope} and
MME \citep{mme}. The base POPE accuracy ($0.877$) matches the published
Janus-Pro-7B value ($0.874$), validating the harness. To avoid certifying ``no
forgetting'' on a target the base already knows, we select each fine-tuning style
by its distance from the base: for text, poetry (Project Gutenberg) is the most
distant candidate at $4.0\times$ the perplexity of a modern-prose control; for
image, a narrow anime style (Naruto, $1{,}141$ pairs) the base renders weakly.
Image-style acquisition is scored against real style images with CMMD
\citep{cmmd} as the robust headline, with KID \citep{kid}, DINOv2
\citep{dinov2}, CLIP-IQA \citep{clipiqa} and CLIP directional similarity
\citep{stylegannada} as complements.

\section{Results}

\subsection{Image-generation direction}
\label{sec:image}
We fine-tune generation to a final configuration (rank $48$, attention and MLP LoRA, a
KL-to-base regularizer
$\mathcal{L}=\mathrm{CE}+\lambda\,\mathrm{KL}(p_{\text{adapter}}\|p_{\text{base}})$,
$\lambda{=}0.5$; the incremental ablation reaching it is in Appendix~\ref{app:recipe}). The adapter
acquires the style (Figure~\ref{fig:teaser}).

\paragraph{No forgetting, modular across capacity.}
Table~\ref{tab:rank} reports the rank sweep against the base. Sweeping the rank from 2 to 64, the adapter acquires the style (CMMD
from $0.139$ to $0.097$, directional similarity from $0.143$ to $0.158$,
saturating around rank 16 to 32), while every forgetting probe stays at the base
level across the whole range. Capacity refines acquisition at no measurable cost
to the other modalities.

\begin{table}[t]
\centering\small
\setlength{\tabcolsep}{6pt}
\begin{tabular}{lccccccc}
\toprule
 & \textbf{Base} & $2$ & $4$ & $8$ & $16$ & $32$ & $64$ \\
\midrule
\multicolumn{8}{l}{\textit{Acquisition (vs.\ real style images)}}\\
CMMD $\downarrow$       & ---     & $.139$ & $.124$ & $.126$ & $.105$ & $.108$ & $\mathbf{.097}$ \\
CLIP DirSim $\uparrow$  & ---     & $.143$ & $.142$ & $.143$ & $.154$ & $.153$ & $\mathbf{.158}$ \\
\midrule
\multicolumn{8}{l}{\textit{Interference}}\\
PPL $\downarrow$        & $13.71$ & $13.66$ & $13.70$ & $13.76$ & $13.74$ & $13.78$ & $13.76$ \\
MMLU $\uparrow$         & $.492$  & $.491$ & $.491$ & $.491$ & $.501$ & $.497$ & $.489$ \\
POPE acc $\uparrow$     & $.877$  & $.879$ & $.881$ & $.878$ & $.878$ & $.879$ & $.878$ \\
MME $\uparrow$          & $1785$  & $1777$ & $1777$ & $1786$ & $1786$ & $1796$ & $1781$ \\
\bottomrule
\end{tabular}
\caption{Image-generation rank sweep against the base. Acquisition (bold = best)
improves and saturates; every forgetting probe stays at the base level at all
ranks. Base POPE ($0.877$) matches the published value ($0.874$).}
\label{tab:rank}
\end{table}

\paragraph{An acquisition optimum in the KL weight.}
At fixed rank $32$, acquisition is a single-peaked function of the KL weight
$\lambda$ (Table~\ref{tab:kl}): it peaks near $\lambda{=}0.25$ to $0.5$ and
collapses for $\lambda\geq1$, where the adapter is over-constrained toward the
base. $\lambda{=}0$ is already workable, so the KL term refines rather than
enables the recipe.

\begin{table}[t]
\centering\small
\setlength{\tabcolsep}{9pt}
\begin{tabular}{lccccc}
\toprule
\textbf{Metric} & $\lambda{=}0$ & $\lambda{=}0.25$ & $\lambda{=}0.5$ & $\lambda{=}1.0$ & $\lambda{=}1.5$ \\
\midrule
CMMD $\downarrow$       & $.155$ & $\mathbf{.100}$ & $.108$ & $.205$ & $.221$ \\
DINO sim $\uparrow$     & $.626$ & $\mathbf{.647}$ & $.643$ & $.439$ & $.388$ \\
CLIP DirSim $\uparrow$  & $.157$ & $\mathbf{.164}$ & $.152$ & $.049$ & $.024$ \\
CLIP-IQA $\uparrow$     & $.666$ & $.691$ & $\mathbf{.691}$ & $.738$ & $.739$ \\

\bottomrule
\end{tabular}
\caption{Image-generation KL sweep at rank $32$. Acquisition (bold = best) peaks
near $\lambda{=}0.25$ to $0.5$ and degrades for $\lambda\geq1$; Appendix~\ref{app:metric-debug}
motivates the $\lambda{=}0.5$ Final recipe from content preservation.}
\label{tab:kl}
\end{table}

\paragraph{Why $\lambda{=}0.5$ for the Final recipe.}
The CLIP style metrics place the acquisition optimum at $\lambda{=}0.25$
(Table~\ref{tab:kl}), yet we adopt $\lambda{=}0.5$. The reason is that CMMD and
directional similarity measure proximity to the style \emph{distribution} but not
whether the prompt's content survives the stylisation, and three content-aware
signals disagree with them: a memorization-robust DINOv2-space distribution
distance and a content-preservation score both peak at $\lambda{=}0.5$
(Appendix~\ref{app:metric-debug}), and an independent VLM judge (Qwen2-VL-7B)
prefers $\lambda{=}0.5$ over $\lambda{=}0.25$ on every criterion
(Table~\ref{tab:judge}), significantly so on content. We mitigate the judge's
position bias by scoring both image orders and counting only order-consistent
preferences (MT-Bench style). The visual intuition behind the Final recipe is thus
confirmed by a rigorous, bias-controlled procedure, while the metric optimum at
$\lambda{=}0.25$ reflects a known limitation of distribution-proximity metrics.

\begin{table}[t]
\centering\small
\setlength{\tabcolsep}{8pt}
\begin{tabular}{lccc}
\toprule
Criterion & $\lambda{=}0.25$ win & $\lambda{=}0.5$ win & tie \\
\midrule
Style    & $23.8\%$ & $\mathbf{30.0\%}$ & $46.2\%$ \\
Content  & $6.2\%$  & $\mathbf{20.0\%}$ & $73.8\%$ \\
Overall  & $23.8\%$ & $\mathbf{30.0\%}$ & $46.2\%$ \\
\bottomrule
\end{tabular}
\caption{Independent VLM judge (Qwen2-VL-7B), $\lambda{=}0.25$ vs.\ $\lambda{=}0.5$
on $80$ held-out prompts, each scored in both image orders. A configuration wins a
prompt only if preferred in \emph{both} orders; order-dependent results are ties
(position-bias mitigation). On the order-consistent comparisons the judge prefers
$\lambda{=}0.5$ on every criterion, significantly on content (decisive win-rate
$76\%$, $95\%$ CI $[55,89]$), and never favours $\lambda{=}0.25$.}
\label{tab:judge}
\end{table}

\paragraph{Beyond a single style: generalisation and a dataset-dependent recipe.}
We replicate the image protocol on a second, unrelated dataset, black-and-white
line-sketch scenes (\texttt{zoheb/sketch-scene}), using the \emph{same}
configuration as Naruto ($r{=}48$, $\alpha{=}96$, $\lambda{=}0.5$, LR $10^{-4}$).
The core result holds (Table~\ref{tab:sketch}): the adapter acquires the style at
the distribution level (CMMD $0.253\!\to\!0.159$, KID $0.263\!\to\!0.096$, DINO
similarity to the sketch reference $0.133\!\to\!0.371$), while the language and
understanding probes show no catastrophic forgetting (WikiText PPL
$13.71\!\to\!14.18$, MMLU $0.492\!\to\!0.479$, POPE $0.877\!\to\!0.871$).
Fine-tuning the generation direction on a completely different image corpus thus
leaves the text and understanding directions intact, replicating the
no-forgetting result of Table~\ref{tab:rank} on a second dataset.

This replication also surfaces a second finding: \emph{the PEFT configuration is
dataset-dependent}. The recipe that saturates a high-frequency, colour-saturated
style (Naruto) only partially imprints the sparse, low-frequency sketch style ~\ref{fig:sketch}.
The per-prompt edit strength is about half (CLIP DirSim $0.086$, versus
$\approx0.15$ for Naruto at the same $\lambda$); it carries a measurable
image-quality cost (CLIP-IQA $0.747\!\to\!0.592$); and the interference
footprint, while small, is uniformly larger than Naruto's and most visible on MME
($1785\!\to\!1713$, which stayed flat for Naruto). A single fixed configuration
therefore does not guarantee equal style acquisition, nor an identical
interference footprint, across datasets; only the qualitative conclusion (no
catastrophic language forgetting) transfers unchanged.

Scaling the recipe does
not help: a stronger configuration (rank $64$, $\lambda{=}0.25$) moves the
aggregate distribution closer to the sketch reference (KID $0.096\!\to\!0.065$,
DINO $0.371\!\to\!0.521$) but does not improve the per-prompt edit strength
(CLIP DirSim $0.086\!\to\!0.089$, still about half of Naruto's) and lowers image
quality (CLIP-IQA $0.59\!\to\!0.49$), indicating the limit is intrinsic to the
low-frequency style rather than a matter of capacity.

\begin{table}[t]
\centering\small
\setlength{\tabcolsep}{12pt}
\begin{tabular}{lcc}
\toprule
\textbf{Metric} & \textbf{Base} & \textbf{Sketch LoRA} \\
\midrule
\multicolumn{3}{l}{\textit{Acquisition (vs.\ sketch reference, $n{=}80$)}}\\
CMMD $\downarrow$        & $.253$  & $\mathbf{.159}_{\pm.007}$ \\
KID $\downarrow$         & $.263$  & $\mathbf{.096}_{\pm.013}$ \\
DINO sim $\uparrow$      & $.133$  & $\mathbf{.371}_{\pm.021}$ \\
CLIP DirSim $\uparrow$   & ---     & $.086_{\pm.006}$ \\
\midrule
\multicolumn{3}{l}{\textit{Interference}}\\
WikiText PPL $\downarrow$ & $13.71_{\pm.21}$ & $14.18_{\pm.22}$ \\
MMLU $\uparrow$           & $.492_{\pm.013}$  & $.479_{\pm.013}$ \\
POPE acc $\uparrow$       & $.877_{\pm.006}$  & $.871_{\pm.006}$ \\
MME $\uparrow$            & $1785_{\pm31}$  & $1713_{\pm31}$ \\
\bottomrule
\end{tabular}
\caption{Generalisation to a second image dataset (line-sketch scenes,
\texttt{zoheb/sketch-scene}) with the \emph{same} configuration as Naruto. The
style is acquired at the distribution level (CMMD/KID/DINO) and language and
knowledge probes show no catastrophic forgetting; the per-prompt edit strength
(CLIP DirSim $0.086$) is about half of Naruto's, evidence that the PEFT recipe is
dataset-dependent. FID/ArtFID are omitted (numerically unstable at $n{=}80$).}
\label{tab:sketch}
\end{table}

\begin{figure}[!t]
\centering
\begin{subfigure}[t]{0.23\textwidth}\centering
\includegraphics[width=\linewidth]{figs/janus_comparison/no_adapter/human_05_natural.png}
\caption*{Base}\end{subfigure}\hfill
\begin{subfigure}[t]{0.23\textwidth}\centering
\includegraphics[width=\linewidth]{figs/janus_comparison/adapter_solo_gen/human_05_natural.png}
\caption*{Image LoRA}\end{subfigure}\hfill
\begin{subfigure}[t]{0.23\textwidth}\centering
\includegraphics[width=\linewidth]{figs/janus_comparison/no_adapter/human_07_natural.png}
\caption*{Base}\end{subfigure}\hfill
\begin{subfigure}[t]{0.23\textwidth}\centering
\includegraphics[width=\linewidth]{figs/janus_comparison/adapter_solo_gen/human_07_natural.png}
\caption*{Image LoRA}\end{subfigure}
\caption{Style acquisition on two held-out prompts. Full grids in
Appendix~\ref{app:grids}.}
\label{fig:teaser}
\end{figure}

\begin{figure}[!t]
\centering
\begin{subfigure}[t]{0.23\textwidth}\centering
\includegraphics[width=\linewidth]{figs/sketch_comparison/zebra_before.png}
\caption*{Base}\end{subfigure}\hfill
\begin{subfigure}[t]{0.23\textwidth}\centering
\includegraphics[width=\linewidth]{figs/sketch_comparison/zebra_after.png}
\caption*{Sketch LoRA}\end{subfigure}\hfill
\begin{subfigure}[t]{0.23\textwidth}\centering
\includegraphics[width=\linewidth]{figs/sketch_comparison/kite_before.png}
\caption*{Base}\end{subfigure}\hfill
\begin{subfigure}[t]{0.23\textwidth}\centering
\includegraphics[width=\linewidth]{figs/sketch_comparison/kite_after.png}
\caption*{Sketch LoRA}\end{subfigure}
\caption{Generalisation to line-sketch scenes (same
prompt/seed base vs.\ adapter). The style transfers for some prompts (left pair:
white canvas, ink linework) but only partially for others (right pair: sketch
strokes appear yet photographic colour persists), consistent with the weaker,
dataset-dependent acquisition in Table~\ref{tab:sketch}.}
\label{fig:sketch}
\end{figure}

\subsection{Text direction}
\label{sec:text}
We now specialize the text language-modelling direction on poetry, the style most
distant from the base by our perplexity probe. We run two sweeps, over the LoRA
rank (Table~\ref{tab:text}) and over the KL weight $\lambda$
(Table~\ref{tab:text-kl}), and for each report task acquisition (held-out poetry
perplexity) together with the forgetting metrics. The text metrics use a
subsampled configuration, so its absolute values are not comparable to the image results of Table~\ref{tab:rank}; the valid comparison is within a table, against
the base column.

\begin{table}[t]
\centering\small
\setlength{\tabcolsep}{7pt}
\begin{tabular}{lcccccc}
\toprule
 & \textbf{Base} & $4$ & $8$ & $16$ & $32$ & $64$ \\
\midrule
\multicolumn{7}{l}{\textit{Acquisition (target style)}}\\
Poetry PPL $\downarrow$   & $32.7$ & $25.6$ & $25.3$ & $24.9$ & $24.4$ & $\mathbf{23.8}$ \\
\midrule
\multicolumn{7}{l}{\textit{Interference}}\\
WikiText PPL $\downarrow$ & $13.71$ & $12.52$ & $12.51$ & $\underline{12.51}$ & $12.53$ & $12.59$ \\
MMLU $\uparrow$           & $.481$ & $.481$ & $.477$ & $.479$ & $.477$ & $.488$ \\
POPE acc $\uparrow$       & $.875$ & $.871$ & $.871$ & $.871$ & $.873$ & $.877$ \\
MME $\uparrow$            & $1130$ & $1129$ & $1111$ & $1128$ & $1130$ & $1122$ \\
\bottomrule
\end{tabular}
\caption{Text (poetry) rank sweep. The style is acquired (poetry perplexity
$32.7$ to $23.8$, bold = best); knowledge and understanding stay flat. WikiText
perplexity is not degraded and is in fact slightly lower at every rank
(underlined best).}
\label{tab:text}
\end{table}

\paragraph{Rank sweep: acquisition without forgetting.}
Increasing the rank from 4 to 64 steadily improves acquisition (poetry perplexity
$32.7$ to $23.8$), while MMLU, POPE and MME stay at the base level at every rank
(Table~\ref{tab:text}). General-text (WikiText) perplexity is not degraded and is
in fact slightly below the base throughout. Specializing the text direction
therefore leaves the other capabilities intact across capacity, mirroring the
image result.
Pushing the rank further to $128$ confirms diminishing returns: acquisition improves
only marginally (poetry perplexity $23.8$ to $23.2$), while a first mild understanding
cost appears (MME $1122$ to $1105$, POPE $0.877$ to $0.867$). We therefore report up to
rank $64$, which sits at the knee of this trade-off.

\paragraph{KL sweep: the regularizer trades acquisition for stability.}
At rank 16 we sweep the KL weight $\lambda$ (Table~\ref{tab:text-kl}). Acquisition
is strongest with little regularization and weakens monotonically as the adapter
is held closer to the base. Knowledge and POPE stay flat; MME shows a small dip at
$\lambda{=}0$, the least-regularized and strongest intervention, which the
regularizer removes. Given the subsampled results and a single seed, we read this
as a mild, suggestive understanding cost at maximal intervention rather than a
conclusive effect.

\begin{table}[t]
\centering\small
\setlength{\tabcolsep}{7pt}
\begin{tabular}{lcccccc}
\toprule
 & \textbf{Base} & $\lambda{=}0$ & $\lambda{=}0.25$ & $\lambda{=}0.5$ & $\lambda{=}1$ & $\lambda{=}1.5$ \\
\midrule
\multicolumn{7}{l}{\textit{Acquisition (target style)}}\\
Poetry PPL $\downarrow$   & $32.7$ & $\mathbf{24.5}$ & $24.5$ & $24.9$ & $25.9$ & $26.7$ \\
\midrule
\multicolumn{7}{l}{\textit{Interference}}\\
WikiText PPL $\downarrow$ & $13.71$ & $12.61$ & $\underline{12.50}$ & $12.51$ & $12.75$ & $12.93$ \\
MMLU $\uparrow$           & $.481$ & $.474$ & $.481$ & $.479$ & $.475$ & $.483$ \\
POPE acc $\uparrow$       & $.875$ & $.868$ & $.870$ & $.871$ & $.878$ & $.873$ \\
MME $\uparrow$            & $1130$ & $1101$ & $1116$ & $1128$ & $1127$ & $1128$ \\
\midrule
\multicolumn{7}{l}{\textit{Generation shift from base (image direction)}}\\
Gen-CMMD $\downarrow$ ($\times10^{-3}$) & $0$ & $1.73$ & $1.71$ & $1.73$ & $1.77$ & $1.65$ \\
\bottomrule
\end{tabular}
\caption{Text (poetry) KL sweep at rank 16. Acquisition (bold = best) is
strongest with little regularization and weakens as $\lambda$ grows. MMLU and
POPE stay flat; MME shows a small dip at $\lambda{=}0$ that the regularizer
removes. WikiText perplexity is below the base at every $\lambda$ (underlined
best). Gen-CMMD is the CMMD between the adapter's and the base's image generations on a
fixed prompt set ($N{=}128$); it stays flat near zero, so the text adapter does not
measurably shift image generation. Subsampled, single seed.}
\label{tab:text-kl}
\end{table}

\paragraph{Beyond likelihood benchmarks.}
The small $\lambda{=}0$ dip is the first benchmark-visible hint of a cost at maximal
intervention, but these benchmarks score only multiple-choice and likelihood items,
never free-form generation. As a direct check on the other generative direction, we
measure the CMMD between the text adapter's and the base's image generations on a
fixed prompt set (Table~\ref{tab:text-kl}, bottom row): it stays flat and near zero
($\sim\!1.7\times10^{-3}$) at every $\lambda$, so specializing the text stream does
not measurably shift image generation. The remaining gap is free-form \emph{text}
generation, which the likelihood results do not capture; we leave its measurement
to future work.


\paragraph{Beyond a single text domain.}
To test that the text-direction result is not specific to poetry, we replicate the
protocol on a second, semantically distant domain, medical abstracts
(\texttt{ccdv/pubmed-summarization}; $r{=}32$, $\lambda{=}0.25$). Acquisition is
substantially stronger than for poetry: held-out PubMed perplexity falls from
$12.8$ to $8.4$, and on the comparable distance probe the domain moves from
$1.23\times$ the general-text control (harder than generic text) to $0.79\times$
(easier than generic text), confirming a large domain shift. Despite this strong
adaptation, the forgetting stays at the base level
(Table~\ref{tab:text-domain}: MMLU $0.492\!\to\!0.495$, POPE $0.877$, MME
$1785\!\to\!1766$, WikiText PPL $13.71\!\to\!13.97$), and the image-generation
direction shifts only marginally: the CMMD between the adapter's and the base's
image generations is $1.77\times10^{-3}$ (Gen-CMMD), the same order as
poetry's $\sim\!1.7$. The central assumption, that specialising the text stream
does not degrade image generation, therefore generalises to a second domain and
holds even under a much stronger textual adaptation than poetry.

\begin{table}[t]
\centering\small
\setlength{\tabcolsep}{12pt}
\begin{tabular}{lcc}
\toprule
\textbf{Metric} & \textbf{Base} & \textbf{Medicine LoRA} \\
\midrule
\multicolumn{3}{l}{\textit{Acquisition (domain text)}}\\
PubMed PPL $\times$control $\downarrow$ & $1.23\times$ & $\mathbf{0.79\times}$ \\
\midrule
\multicolumn{3}{l}{\textit{Interference}}\\
WikiText PPL $\downarrow$ & $13.71_{\pm.21}$ & $13.97_{\pm.21}$ \\
MMLU $\uparrow$           & $.492_{\pm.013}$  & $.495_{\pm.013}$ \\
POPE acc $\uparrow$       & $.877_{\pm.006}$  & $.877_{\pm.006}$ \\
MME $\uparrow$            & $1785_{\pm31}$  & $1766_{\pm32}$ \\
\midrule
\multicolumn{3}{l}{\textit{Generation shift from base (image direction)}}\\
Gen-CMMD $\downarrow$ ($\times10^{-3}$) & $0$ & $1.77$ \\
\bottomrule
\end{tabular}
\caption{Generalisation of the text direction to a second, distant domain
(medicine, \texttt{ccdv/pubmed-summarization}; $r{=}32$, $\lambda{=}0.25$). The
adapter strongly acquires the domain (PubMed perplexity from $1.23\times$ to
$0.79\times$ the general-text control, i.e.\ medical text becomes \emph{easier}
than generic text), while the forgetting stays at the base level and image
generation is essentially unchanged (Gen-CMMD near poetry's). Strong textual
adaptation therefore does not degrade the image-generation direction.
This table uses the \emph{full} forgetting probes, so its base absolute values
(e.g.\ MME $1785$) differ from the light-config poetry/Welsh tables; the
comparison is within-table against the base column.}
\label{tab:text-domain}
\end{table}

\paragraph{A genuinely unseen target: a new language.}
Poetry and medicine both lie within the base's competence (poetic English,
medical English), so the absence of forgetting could be attributed to a weak
intervention. We therefore add a third, deliberately harder target: a different
\emph{language}, Welsh (\texttt{wikimedia/wikipedia}, Welsh subset; $r{=}32$,
$\lambda{=}0.1$), selected by a distance probe over candidate low-resource and
constructed languages. Acquisition is the strongest of the three text targets:
held-out Welsh perplexity falls from $40.7$ to $14.6$ ($-64\%$), and on Welsh
seeds the adapter sustains Welsh where the base slips into English or its
assistant persona (Appendix~\ref{app:welsh-examples}). Despite acquiring an entire
new language, the forgetting stays at the base level (Table~\ref{tab:welsh}: MMLU
$.481\!\to\!.479$, POPE $.875\!\to\!.876$, MME $1130\!\to\!1124$, WikiText PPL
$13.71\!\to\!12.07$), and the image-generation direction is essentially unchanged:
the CMMD between the adapter's and the base's generations is $1.7\times10^{-3}$
(Gen-CMMD), the same order as poetry and medicine. Two observations sharpen the
result. First, this strongest acquisition requires the \emph{smallest} weight
change of the three text adapters ($\lVert\Delta W\rVert/\lVert W\rVert = 0.6\%$,
vs.\ $1.3\%$ for poetry and $1.7\%$ for medicine), despite using the weakest
regularizer ($\lambda{=}0.1$): the new capability appears to be unlocked rather
than built, i.e.\ partly latent in the pretrained backbone. Second, under the
adapter English prompts are still continued in English (no language bleed), so the
language is added without overwriting the existing one. The central assumption ---
that specializing the text stream does not degrade image generation --- therefore
holds even under the acquisition of a new language.

\begin{table}[t]
\centering\small
\setlength{\tabcolsep}{12pt}
\begin{tabular}{lcc}
\toprule
\textbf{Metric} & \textbf{Base} & \textbf{Welsh LoRA} \\
\midrule
\multicolumn{3}{l}{\textit{Acquisition (target language)}}\\
Welsh held-out PPL $\downarrow$ & $40.7$ & $\mathbf{14.6}$ \\
$\lVert\Delta W\rVert/\lVert W\rVert$ & --- & $0.6\%$ \\
\midrule
\multicolumn{3}{l}{\textit{Interference}}\\
WikiText PPL $\downarrow$ & $13.71_{\pm.21}$ & $12.07_{\pm.18}$ \\
MMLU $\uparrow$           & $.481_{\pm.021}$  & $.479_{\pm.016}$ \\
POPE acc $\uparrow$       & $.875_{\pm.010}$  & $.876_{\pm.010}$ \\
MME $\uparrow$            & $1130_{\pm42}$  & $1124_{\pm42}$ \\\midrule
\multicolumn{3}{l}{\textit{Generation shift from base (image direction)}}\\
Gen-CMMD $\downarrow$ ($\times10^{-3}$) & $0$ & $1.73$ \\
\bottomrule
\end{tabular}
\caption{Welsh language target (\texttt{wikimedia/wikipedia}, Welsh; $r{=}32$,
$\lambda{=}0.1$). The adapter strongly acquires the language (held-out perplexity
$40.7\!\to\!14.6$, $-64\%$) --- the largest acquisition of the three text
targets --- yet needs the smallest weight change ($0.6\%$). Knowledge (MMLU),
grounded understanding (POPE/MME) and general English (WikiText) stay at the base
level, and image generation is essentially unchanged (Gen-CMMD near
poetry/medicine). Light probe config (base column = base\_light); standard
errors from evaluation-set resampling.}
\label{tab:welsh}
\end{table}

\section{Discussion}
\paragraph{Where modularity holds.}
On standard, likelihood-scored benchmarks, PEFT on Janus-Pro-7B is modular:
specializing image generation leaves text and understanding unchanged across rank
and KL (Section~\ref{sec:image}), and specializing text preserves knowledge and
understanding while acquiring the style (Section~\ref{sec:text}). The open
question is whether the same holds for free-form generative quality at high
capacity, which the present benchmarks do not test.

\paragraph{How large is the intervention?}
Is the adapter's footprint large enough that ``no forgetting'' is non-trivial? We
summarize each update $\Delta W=(\alpha/r)BA$ by its relative Frobenius norm
$\lVert\Delta W\rVert_F/\lVert W\rVert_F$, a task-vector magnitude
\citep{ilharco2023task}. It is small but real, $0.6$--$1.9\%$ of the backbone's
weight energy, and tracks the two knobs we control in both directions: rising
with rank and shrinking as the KL weight grows (e.g.\ at $r{=}32$, image $1.6\%$
vs.\ text $1.3\%$; Appendix Tables~\ref{tab:dw-rank},~\ref{tab:dw-kl}). Yet these
sub-$2\%$ updates already shift behavior (poetry perplexity $32.7{\to}23.8$;
visible image-style transfer) while the other modality stays flat: a targeted,
low-norm, modality-localized perturbation suffices, the modularity we claim. The
heatmaps (Appendix Fig.~\ref{fig:dw-heatmap}) confirm it is localized, attention
$q/k/v$ near-untouched, the change concentrated in the MLP (broad \texttt{gate}
for text, final-layer \texttt{down} for image), so the global percentage
understates the local effect.
\section{Conclusion and Future Work}
We studied the modularity of LoRA fine-tuning on a unified multimodal model. With
fine-tuning targets chosen by distance from the base, specializing the generation
direction leaves the other capabilities unchanged on public benchmarks, robustly
across rank and KL, and specializing the text direction preserves knowledge and
understanding while acquiring the style. Future work: quantifying free-form
generative quality with distributional metrics (MAUVE for text; CMMD and
FID/ArtFID at scale for images) and averaging over multiple training seeds, which
would also let us test the small WikiText-perplexity improvement as genuine
positive transfer; the adapter-placement study along Janus' decoupled pathways;
replication on additional styles (pixel art); and extending to joint
multi-direction training and to a third direction (visual question answering).

\bibliographystyle{plainnat}
\bibliography{references}

\appendix

\section{Implementation details}
\label{app:impl}

\paragraph{Model.}
All experiments use Janus-Pro-7B (\texttt{deepseek-community/Janus-Pro-7B}),
loaded in \texttt{bfloat16} via the Hugging Face
\texttt{JanusForConditionalGeneration} class. Janus-Pro is a unified
autoregressive model with a shared transformer backbone and \emph{decoupled}
modality encoders: a SigLIP vision encoder for understanding (image$\to$text) and
a separate vector-quantized (VQ) tokenizer with an image head for generation
(text$\to$image), where an image is a $24\times24=576$ grid of discrete tokens
from a codebook of size $16384$. We use the two directions through the
processor's text mode (text language modelling) and image mode
(\texttt{max\_new\_tokens}$=576$ then VQ decoding). LoRA is applied \emph{only} to
the shared backbone; the SigLIP and VQ encoders/heads are frozen, which is what
makes the modality-localized weight analysis of Appendix~\ref{app:weight-delta}
meaningful.

\paragraph{LoRA configuration (fixed across all runs).}
Every adapter targets the seven projections of each backbone layer---%
\texttt{q,k,v,o\_proj} (attention) and \texttt{gate,up,down\_proj} (MLP)---with
$\alpha=2r$ (so the LoRA scaling $\alpha/r=2$ is constant across the rank sweep),
\texttt{bias=none}, causal-LM task type, and no rsLoRA; the induced update is
$\Delta W=(\alpha/r)BA$. Only the swept variable (rank $r$ or KL weight $\lambda$)
changes within a sweep.

\paragraph{Optimization (common).}
AdamW with $\beta=(0.9,0.95)$ (matching the Janus pre-training $\beta_2$) and zero
weight decay, a cosine schedule with $5\%$ warmup, gradient accumulation $8$,
gradient clipping $1.0$, \texttt{bfloat16}, single H100. With the regularizer the
loss is
$\mathcal{L}=\mathrm{CE}+\lambda\,\mathrm{KL}(p_{\text{adapter}}\|p_{\text{base}})$,
where $p_{\text{base}}$ is a \texttt{disable\_adapter()} no-grad forward on the
same inputs ($\approx$$2\times$ forward cost).

\paragraph{Uncertainty quantification.}
We report each number as $\text{mean}\pm$ one standard error, computed without
retraining and with an estimator-appropriate method per metric type. The accuracy
benchmarks (MMLU, POPE, MME) are proportions of correct items over a fixed
evaluation set, so their standard error follows the binomial/Wilson form
\citep{wilson1927,brown2001interval}. Per-instance metrics averaged over the
generated set (CLIP directional similarity, DINO similarity, CLIP-IQA) use
$s/\sqrt{n}$ over the $n$ generations, equivalently a nonparametric bootstrap
\citep{efron1993bootstrap}. The set-level distributional metrics (CMMD, KID) have
no closed-form variance, so their standard error is the standard deviation of a
nonparametric bootstrap over the generated and reference sets ($B{=}2000$
resamples) \citep{efron1993bootstrap,binkowski2018demystifying}. Perplexity is a
monotone transform of the mean per-chunk negative log-likelihood, whose standard
error propagates through the exponential. All standard errors quantify
evaluation-set sampling noise at fixed model weights; they do not capture
training-seed variance. The small benchmark movements we observe (e.g.\ MMLU
$0.492\!\to\!0.479$, within $\sim\!1$ standard error) thus lie within sampling
noise, whereas the acquisition shifts (e.g.\ CMMD, DirSim) do not.

\paragraph{Text (poetry) runs.}
Data: \texttt{biglam/gutenberg-poetry-corpus}, grouped by source document with
front-matter removed and windowed to $512$ tokens ($6000$ train / $200$ val
windows). Fixed: $r{=}16$, $\alpha{=}32$, dropout $0.05$, LR $10^{-4}$, $3$
epochs ($150$-min wall-clock budget). Rank sweep: $r\in\{4,8,16,32,64\}$ at
$\lambda{=}0.5$; $\lambda$ sweep: $\lambda\in\{0,0.25,0.5,1,1.5\}$ at $r{=}16$.

\paragraph{Image (Naruto) runs.}
Data: \texttt{lambdalabs/naruto-blip-captions} ($\approx$$1141$ image--caption
pairs); acquisition on $80$ held-out prompts. Final configuration
(Table~\ref{tab:panel}): $r{=}48$, $\alpha{=}96$, dropout $0.10$, LR $10^{-4}$,
$6$ epochs, no caption prefix, $\lambda{=}0.5$. Rank sweep:
$r\in\{2,\dots,64\}$ at $\lambda{=}0.5$; $\lambda$ sweep at $r{=}32$.

\paragraph{Prompts and image generation.}
Janus generates images autoregressively from a natural-language prompt
(text$\to$image); every generated image in this paper is therefore conditioned on
an explicit text prompt. During training the conditioning text is the dataset's
BLIP caption; the \textit{Capacity} and \textit{+KL} configurations prepend a
fixed style trigger (``anime sprite in Naruto style, ''), which \textit{Scaled}
and \textit{Final} remove, so the Final adapter applies the style
\emph{unconditionally}, with no trigger token at inference. Acquisition metrics
(Table~\ref{tab:panel}) are computed on images generated from $80$ held-out Naruto
captions. The qualitative grids instead use a fixed, curated set of generic
natural-language prompts mixing in-domain and clearly out-of-domain subjects
(e.g.\ ``a young man with blonde spiky hair'', ``a samurai warrior holding a
katana'', ``an old wizard with a long white beard''), chosen to test whether the
style transfers to arbitrary content; the exact prompt is printed beneath each row
of every grid (Figures~\ref{fig:janus_qualitative_comparison},
\ref{fig:rank_visual}, \ref{fig:kl_visual}). The cross-modal check under the text
adapter (Figure~\ref{fig:t2i_text_kl}) uses four neutral prompts (e.g.\ ``a cat
sitting on a chair''). All conditions share the same seed schedule, so within a
row the adapter is the only source of variation.


\section{Image-generation configuration: incremental ablation}
\label{app:recipe}
We reach the Final configuration by adding components one or a few at a time from
a naive starting point (attention-only LoRA, rank 8, with a fixed caption-prefix
trigger), which degraded the base and is omitted from Table~\ref{tab:panel}.
\textit{Capacity} restores trainable capacity (rank 16, LoRA on attention and
MLP); \textit{+KL} adds the KL-to-base regularizer ($\lambda{=}0.5$);
\textit{Scaled} removes the training-time trigger prefix and raises rank, learning
rate and dropout (rank 32, $8\times10^{-5}$, dropout 0.1, six epochs);
\textit{Final} raises the rank to 48 and the learning rate to $10^{-4}$. The
\textit{Joint} conditions additionally alternate a $50/50$ domain-controlled text
stream (matched: Naruto franchise; far: physics and mathematics) into the same
adapter. Table~\ref{tab:panel} reports the acquisition panel; the corresponding
generations are in Figure~\ref{fig:janus_qualitative_comparison}.

\begin{table}[t]
\centering\small
\setlength{\tabcolsep}{5pt}
\renewcommand{\arraystretch}{1.05}
\begin{tabular}{lcccccc}
\toprule
 & \textbf{Capacity} & \textbf{+KL} & \textbf{Scaled} & \textbf{Final} & \textbf{Joint} & \textbf{Joint} \\
 & & & & & \textbf{(matched)} & \textbf{(far)} \\
\midrule
\multicolumn{7}{l}{\textit{Configuration (bold = changed vs.\ previous column)}}\\
Rank $r$              & 16 & 16 & \textbf{32} & \textbf{48} & 48 & 48 \\
LR ($\times10^{-5}$)  & 5 & 5 & \textbf{8} & \textbf{10} & 10 & 10 \\
Epochs               & 8 & 8 & \textbf{6} & 6 & 6 & 6 \\
LoRA dropout         & .05 & .05 & \textbf{.10} & .10 & .10 & .10 \\
Train prefix         & yes & yes & \textbf{no} & no & no & no \\
KL $\lambda$         & 0 & \textbf{0.5} & 0.5 & 0.5 & 0.5 & 0.5 \\
Joint text           & --- & --- & --- & --- & \textbf{matched} & \textbf{far} \\
\midrule
\multicolumn{7}{l}{\textit{Acquisition (vs.\ real style, $N{=}80$; bold = best)}}\\
KID $\downarrow$       & $.039$ & $.017$ & $.016$ & $.016$ & $.016$ & $\mathbf{.014}$ \\
CMMD $\downarrow$      & $.189$ & $\mathbf{.105}$ & $.120$ & $.110$ & $.131$ & $.134$ \\
DINO sim $\uparrow$    & $.629$ & $\mathbf{.663}$ & $.624$ & $.640$ & $.633$ & $.637$ \\
CLIP-IQA $\uparrow$    & $.659$ & $\mathbf{.750}$ & $.694$ & $.701$ & $.725$ & $.704$ \\
CLIP DirSim $\uparrow$ & $-.059$ & $-.033$ & $.157$ & $\mathbf{.163}$ & $.146$ & $.151$ \\
\bottomrule
\end{tabular}
\caption{Image-generation configuration as an incremental ablation: each column
adds one or more components to the previous one to reach \textit{Final}.
\textit{Capacity}: rank 16 with LoRA on attention and MLP. \textit{+KL}: adds the
KL-to-base regularizer. \textit{Scaled}: removes the training-time caption prefix
and raises rank, learning rate and dropout. \textit{Final}: rank 48, LR
$10^{-4}$. \textit{Joint}: alternates a $50/50$ domain-controlled text stream
(matched: Naruto franchise; far: physics/mathematics) into the same adapter. The
naive start (attention-only, rank 8, prefix-gated) degraded the base and is
omitted; CLIP DirSim turning positive at \textit{Scaled} marks the onset of
genuine style transfer.}
\label{tab:panel}
\end{table}

\begin{figure*}[p]
\centering
\setlength{\tabcolsep}{1.5pt}
\renewcommand{\arraystretch}{1.0}
\small
\begin{tabular}{ccccccc}
\textbf{No adapter} & \textbf{Capacity} & \textbf{+KL} & \textbf{Scaled} & \textbf{Final} & \textbf{Joint (matched)} & \textbf{Joint (far)} \\
\textit{base 7B} & \textit{r16, +MLP} & \textit{+ KL reg} & \textit{r32, untriggered} & \textit{r48, LR $10^{-4}$} & \textit{+ Naruto wiki} & \textit{+ physics wiki} \\
\midrule
\includegraphics[width=0.135\textwidth]{figs/janus_comparison/no_adapter/eval_00_natural.png} &
\includegraphics[width=0.135\textwidth]{figs/janus_comparison/adapter_v2/eval_00_natural.png} &
\includegraphics[width=0.135\textwidth]{figs/janus_comparison/adapter_v3/eval_00_natural.png} &
\includegraphics[width=0.135\textwidth]{figs/janus_comparison/adapter_v4/eval_00_natural.png} &
\includegraphics[width=0.135\textwidth]{figs/janus_comparison/adapter_solo_gen/eval_00_natural.png} &
\includegraphics[width=0.135\textwidth]{figs/janus_comparison/adapter_text_matched/eval_00_natural.png} &
\includegraphics[width=0.135\textwidth]{figs/janus_comparison/adapter_text_far/eval_00_natural.png} \\
\multicolumn{7}{l}{\footnotesize\textit{Prompt (eval\_00):} a red panda with a black face and brown eyes} \\
\addlinespace[2pt]
\includegraphics[width=0.135\textwidth]{figs/janus_comparison/no_adapter/eval_01_natural.png} &
\includegraphics[width=0.135\textwidth]{figs/janus_comparison/adapter_v2/eval_01_natural.png} &
\includegraphics[width=0.135\textwidth]{figs/janus_comparison/adapter_v3/eval_01_natural.png} &
\includegraphics[width=0.135\textwidth]{figs/janus_comparison/adapter_v4/eval_01_natural.png} &
\includegraphics[width=0.135\textwidth]{figs/janus_comparison/adapter_solo_gen/eval_01_natural.png} &
\includegraphics[width=0.135\textwidth]{figs/janus_comparison/adapter_text_matched/eval_01_natural.png} &
\includegraphics[width=0.135\textwidth]{figs/janus_comparison/adapter_text_far/eval_01_natural.png} \\
\multicolumn{7}{l}{\footnotesize\textit{Prompt (eval\_01):} a man with glasses and a hoodie} \\
\addlinespace[2pt]
\includegraphics[width=0.135\textwidth]{figs/janus_comparison/no_adapter/eval_02_natural.png} &
\includegraphics[width=0.135\textwidth]{figs/janus_comparison/adapter_v2/eval_02_natural.png} &
\includegraphics[width=0.135\textwidth]{figs/janus_comparison/adapter_v3/eval_02_natural.png} &
\includegraphics[width=0.135\textwidth]{figs/janus_comparison/adapter_v4/eval_02_natural.png} &
\includegraphics[width=0.135\textwidth]{figs/janus_comparison/adapter_solo_gen/eval_02_natural.png} &
\includegraphics[width=0.135\textwidth]{figs/janus_comparison/adapter_text_matched/eval_02_natural.png} &
\includegraphics[width=0.135\textwidth]{figs/janus_comparison/adapter_text_far/eval_02_natural.png} \\
\multicolumn{7}{l}{\footnotesize\textit{Prompt (eval\_02):} a person with black hair and a green shirt} \\
\addlinespace[2pt]
\includegraphics[width=0.135\textwidth]{figs/janus_comparison/no_adapter/human_05_natural.png} &
\includegraphics[width=0.135\textwidth]{figs/janus_comparison/adapter_v2/human_05_natural.png} &
\includegraphics[width=0.135\textwidth]{figs/janus_comparison/adapter_v3/human_05_natural.png} &
\includegraphics[width=0.135\textwidth]{figs/janus_comparison/adapter_v4/human_05_natural.png} &
\includegraphics[width=0.135\textwidth]{figs/janus_comparison/adapter_solo_gen/human_05_natural.png} &
\includegraphics[width=0.135\textwidth]{figs/janus_comparison/adapter_text_matched/human_05_natural.png} &
\includegraphics[width=0.135\textwidth]{figs/janus_comparison/adapter_text_far/human_05_natural.png} \\
\multicolumn{7}{l}{\footnotesize\textit{Prompt (human\_05):} a young man with blonde spiky hair} \\
\addlinespace[2pt]
\includegraphics[width=0.135\textwidth]{figs/janus_comparison/no_adapter/human_06_natural.png} &
\includegraphics[width=0.135\textwidth]{figs/janus_comparison/adapter_v2/human_06_natural.png} &
\includegraphics[width=0.135\textwidth]{figs/janus_comparison/adapter_v3/human_06_natural.png} &
\includegraphics[width=0.135\textwidth]{figs/janus_comparison/adapter_v4/human_06_natural.png} &
\includegraphics[width=0.135\textwidth]{figs/janus_comparison/adapter_solo_gen/human_06_natural.png} &
\includegraphics[width=0.135\textwidth]{figs/janus_comparison/adapter_text_matched/human_06_natural.png} &
\includegraphics[width=0.135\textwidth]{figs/janus_comparison/adapter_text_far/human_06_natural.png} \\
\multicolumn{7}{l}{\footnotesize\textit{Prompt (human\_06):} a samurai warrior holding a katana} \\
\addlinespace[2pt]
\includegraphics[width=0.135\textwidth]{figs/janus_comparison/no_adapter/human_07_natural.png} &
\includegraphics[width=0.135\textwidth]{figs/janus_comparison/adapter_v2/human_07_natural.png} &
\includegraphics[width=0.135\textwidth]{figs/janus_comparison/adapter_v3/human_07_natural.png} &
\includegraphics[width=0.135\textwidth]{figs/janus_comparison/adapter_v4/human_07_natural.png} &
\includegraphics[width=0.135\textwidth]{figs/janus_comparison/adapter_solo_gen/human_07_natural.png} &
\includegraphics[width=0.135\textwidth]{figs/janus_comparison/adapter_text_matched/human_07_natural.png} &
\includegraphics[width=0.135\textwidth]{figs/janus_comparison/adapter_text_far/human_07_natural.png} \\
\multicolumn{7}{l}{\footnotesize\textit{Prompt (human\_07):} an old wizard with a long white beard} \\
\addlinespace[2pt]
\includegraphics[width=0.135\textwidth]{figs/janus_comparison/no_adapter/human_08_natural.png} &
\includegraphics[width=0.135\textwidth]{figs/janus_comparison/adapter_v2/human_08_natural.png} &
\includegraphics[width=0.135\textwidth]{figs/janus_comparison/adapter_v3/human_08_natural.png} &
\includegraphics[width=0.135\textwidth]{figs/janus_comparison/adapter_v4/human_08_natural.png} &
\includegraphics[width=0.135\textwidth]{figs/janus_comparison/adapter_solo_gen/human_08_natural.png} &
\includegraphics[width=0.135\textwidth]{figs/janus_comparison/adapter_text_matched/human_08_natural.png} &
\includegraphics[width=0.135\textwidth]{figs/janus_comparison/adapter_text_far/human_08_natural.png} \\
\multicolumn{7}{l}{\footnotesize\textit{Prompt (human\_08):} a teenage girl with pink hair} \\
\end{tabular}
\caption{Qualitative comparison across the seven training conditions on
Janus-Pro-7B. Columns 1--5 trace the image-only LoRA progression to the
\textit{Final} configuration; columns 6--7 add a joint text-training stream from
two semantically distant Wikipedia domains (Naruto franchise vs.\
physics/mathematics). All inferences share the same seed schedule; the adapter is
the only source of variation within a row.}
\label{fig:janus_qualitative_comparison}
\end{figure*}

\paragraph{Joint training: a corpus-size effect on perplexity.}
Adding a text stream inflates held-out text perplexity only when the text corpus
is small (Table~\ref{tab:joint}). The controlled far@320 run (distant domain,
matched chunk budget) reaches $55.1$, higher than the near-domain matched run,
while MMLU and POPE stay flat. The inflation is therefore a corpus-size effect (a
small corpus is revisited and over-fit), not a domain-distance effect, and it
does not correspond to any degradation in the other capability directions.

\begin{table}[t]
\centering\small
\setlength{\tabcolsep}{6pt}
\begin{tabular}{llcccc}
\toprule
\textbf{Condition} & \textbf{Text corpus} & \textbf{PPL}$\downarrow$ & \textbf{MMLU}$\uparrow$ & \textbf{POPE acc}$\uparrow$ & \textbf{MME}$\uparrow$ \\
\midrule
Base (no adapter) & ---                  & $13.71$ & $0.492$ & $0.877$ & $1785$ \\
Final (image-only)& ---                  & $13.80$ & $0.491$ & $0.879$ & $1789$ \\
Joint matched     & near, 320 chunks     & $47.87$ & $0.505$ & $0.876$ & $1747$ \\
Joint far         & distant, 7672 chunks & $12.46$ & $0.491$ & $0.878$ & $1795$ \\
Joint far@320     & distant, 320 chunks  & $55.10$ & $0.481$ & $0.873$ & $1736$ \\
\bottomrule
\end{tabular}
\caption{Corpus-size control on the forgetting. Perplexity inflation
tracks corpus size, not domain distance; knowledge and understanding are
unaffected.}
\label{tab:joint}
\end{table}


\section{Qualitative grids}
\label{app:grids}

\begin{figure*}[p]
\centering
\setlength{\tabcolsep}{1.5pt}\small
\renewcommand{\arraystretch}{1.0}
\begin{tabular}{ccccccc}
\textbf{Base} & \textbf{r=2} & \textbf{r=4} & \textbf{r=8} & \textbf{r=16} & \textbf{r=32} & \textbf{r=64} \\
\midrule
\includegraphics[width=0.135\textwidth]{figs/rank_comparison/base/eval_00_natural.png} &
\includegraphics[width=0.135\textwidth]{figs/rank_comparison/r2/eval_00_natural.png} &
\includegraphics[width=0.135\textwidth]{figs/rank_comparison/r4/eval_00_natural.png} &
\includegraphics[width=0.135\textwidth]{figs/rank_comparison/r8/eval_00_natural.png} &
\includegraphics[width=0.135\textwidth]{figs/rank_comparison/r16/eval_00_natural.png} &
\includegraphics[width=0.135\textwidth]{figs/rank_comparison/r32/eval_00_natural.png} &
\includegraphics[width=0.135\textwidth]{figs/rank_comparison/r64/eval_00_natural.png} \\
\multicolumn{7}{l}{\footnotesize\textit{Prompt:} a red panda with a black face and brown eyes} \\
\addlinespace[2pt]
\includegraphics[width=0.135\textwidth]{figs/rank_comparison/base/eval_01_natural.png} &
\includegraphics[width=0.135\textwidth]{figs/rank_comparison/r2/eval_01_natural.png} &
\includegraphics[width=0.135\textwidth]{figs/rank_comparison/r4/eval_01_natural.png} &
\includegraphics[width=0.135\textwidth]{figs/rank_comparison/r8/eval_01_natural.png} &
\includegraphics[width=0.135\textwidth]{figs/rank_comparison/r16/eval_01_natural.png} &
\includegraphics[width=0.135\textwidth]{figs/rank_comparison/r32/eval_01_natural.png} &
\includegraphics[width=0.135\textwidth]{figs/rank_comparison/r64/eval_01_natural.png} \\
\multicolumn{7}{l}{\footnotesize\textit{Prompt:} a man with glasses and a hoodie} \\
\addlinespace[2pt]
\includegraphics[width=0.135\textwidth]{figs/rank_comparison/base/human_05_natural.png} &
\includegraphics[width=0.135\textwidth]{figs/rank_comparison/r2/human_05_natural.png} &
\includegraphics[width=0.135\textwidth]{figs/rank_comparison/r4/human_05_natural.png} &
\includegraphics[width=0.135\textwidth]{figs/rank_comparison/r8/human_05_natural.png} &
\includegraphics[width=0.135\textwidth]{figs/rank_comparison/r16/human_05_natural.png} &
\includegraphics[width=0.135\textwidth]{figs/rank_comparison/r32/human_05_natural.png} &
\includegraphics[width=0.135\textwidth]{figs/rank_comparison/r64/human_05_natural.png} \\
\multicolumn{7}{l}{\footnotesize\textit{Prompt:} a young man with blonde spiky hair} \\
\addlinespace[2pt]
\includegraphics[width=0.135\textwidth]{figs/rank_comparison/base/human_07_natural.png} &
\includegraphics[width=0.135\textwidth]{figs/rank_comparison/r2/human_07_natural.png} &
\includegraphics[width=0.135\textwidth]{figs/rank_comparison/r4/human_07_natural.png} &
\includegraphics[width=0.135\textwidth]{figs/rank_comparison/r8/human_07_natural.png} &
\includegraphics[width=0.135\textwidth]{figs/rank_comparison/r16/human_07_natural.png} &
\includegraphics[width=0.135\textwidth]{figs/rank_comparison/r32/human_07_natural.png} &
\includegraphics[width=0.135\textwidth]{figs/rank_comparison/r64/human_07_natural.png} \\
\multicolumn{7}{l}{\footnotesize\textit{Prompt:} an old wizard with a long white beard} \\
\addlinespace[2pt]
\includegraphics[width=0.135\textwidth]{figs/rank_comparison/base/human_08_natural.png} &
\includegraphics[width=0.135\textwidth]{figs/rank_comparison/r2/human_08_natural.png} &
\includegraphics[width=0.135\textwidth]{figs/rank_comparison/r4/human_08_natural.png} &
\includegraphics[width=0.135\textwidth]{figs/rank_comparison/r8/human_08_natural.png} &
\includegraphics[width=0.135\textwidth]{figs/rank_comparison/r16/human_08_natural.png} &
\includegraphics[width=0.135\textwidth]{figs/rank_comparison/r32/human_08_natural.png} &
\includegraphics[width=0.135\textwidth]{figs/rank_comparison/r64/human_08_natural.png} \\
\multicolumn{7}{l}{\footnotesize\textit{Prompt:} a teenage girl with pink hair} \\
\end{tabular}
\caption{Qualitative rank sweep (Final configuration): the target style is already
present at $r{=}2$ and only refines with rank, remaining stable across the range.}
\label{fig:rank_visual}
\end{figure*}

\begin{figure}[t]
\centering
\includegraphics[width=\textwidth]{figs/rank_sweep/rank_sweep_figure.pdf}
\caption{Rank sweep on the Final configuration (absolute metric values).
\textbf{Left:} visual-style acquisition (CLIP directional similarity); the
dotted line at $0$ marks an adapter with no directional effect. Acquisition is
positive at every rank and rises modestly, saturating around $r=16$--$32$.
\textbf{Right:} visual understanding (POPE accuracy); the dashed line is the
base model. The adapter tracks the base across the whole range, i.e.\ no
cross-modal interference at any rank.}
\label{fig:rank_sweep_panels}
\end{figure}

\begin{figure}[t]
\centering
\includegraphics[width=0.7\columnwidth]{figs/kl_sweep/kl_sweep_image.pdf}
\caption{Style acquisition (CLIP directional similarity, absolute scale)
versus the KL weight $\lambda$ at rank $32$. The dotted line at $0$ marks an
adapter with no directional effect. Acquisition peaks at
$\lambda\!\approx\!0.25$--$0.5$ and collapses toward $0$ for $\lambda\!\geq\!1$,
where the adapter is over-regularised toward the base.}
\label{fig:kl_curve}
\end{figure}

\begin{figure*}[p]
\centering
\setlength{\tabcolsep}{1.5pt}\small
\begin{tabular}{cccccc}
\textbf{Base} & \textbf{$\lambda{=}0$} & \textbf{$\lambda{=}0.25$} & \textbf{$\lambda{=}0.5$} & \textbf{$\lambda{=}1.0$} & \textbf{$\lambda{=}1.5$} \\
\midrule
\includegraphics[width=0.155\textwidth]{figs/kl_comparison/base/eval_00_natural.png} &
\includegraphics[width=0.155\textwidth]{figs/kl_comparison/kl0/eval_00_natural.png} &
\includegraphics[width=0.155\textwidth]{figs/kl_comparison/kl025/eval_00_natural.png} &
\includegraphics[width=0.155\textwidth]{figs/kl_comparison/kl05/eval_00_natural.png} &
\includegraphics[width=0.155\textwidth]{figs/kl_comparison/kl1/eval_00_natural.png} &
\includegraphics[width=0.155\textwidth]{figs/kl_comparison/kl15/eval_00_natural.png} \\
\multicolumn{6}{l}{\footnotesize\textit{Prompt:} a red panda with a black face and brown eyes} \\
\addlinespace[2pt]
\includegraphics[width=0.155\textwidth]{figs/kl_comparison/base/eval_01_natural.png} &
\includegraphics[width=0.155\textwidth]{figs/kl_comparison/kl0/eval_01_natural.png} &
\includegraphics[width=0.155\textwidth]{figs/kl_comparison/kl025/eval_01_natural.png} &
\includegraphics[width=0.155\textwidth]{figs/kl_comparison/kl05/eval_01_natural.png} &
\includegraphics[width=0.155\textwidth]{figs/kl_comparison/kl1/eval_01_natural.png} &
\includegraphics[width=0.155\textwidth]{figs/kl_comparison/kl15/eval_01_natural.png} \\
\multicolumn{6}{l}{\footnotesize\textit{Prompt:} a man with glasses and a hoodie} \\
\addlinespace[2pt]
\includegraphics[width=0.155\textwidth]{figs/kl_comparison/base/human_05_natural.png} &
\includegraphics[width=0.155\textwidth]{figs/kl_comparison/kl0/human_05_natural.png} &
\includegraphics[width=0.155\textwidth]{figs/kl_comparison/kl025/human_05_natural.png} &
\includegraphics[width=0.155\textwidth]{figs/kl_comparison/kl05/human_05_natural.png} &
\includegraphics[width=0.155\textwidth]{figs/kl_comparison/kl1/human_05_natural.png} &
\includegraphics[width=0.155\textwidth]{figs/kl_comparison/kl15/human_05_natural.png} \\
\multicolumn{6}{l}{\footnotesize\textit{Prompt:} a young man with blonde spiky hair} \\
\addlinespace[2pt]
\includegraphics[width=0.155\textwidth]{figs/kl_comparison/base/human_07_natural.png} &
\includegraphics[width=0.155\textwidth]{figs/kl_comparison/kl0/human_07_natural.png} &
\includegraphics[width=0.155\textwidth]{figs/kl_comparison/kl025/human_07_natural.png} &
\includegraphics[width=0.155\textwidth]{figs/kl_comparison/kl05/human_07_natural.png} &
\includegraphics[width=0.155\textwidth]{figs/kl_comparison/kl1/human_07_natural.png} &
\includegraphics[width=0.155\textwidth]{figs/kl_comparison/kl15/human_07_natural.png} \\
\multicolumn{6}{l}{\footnotesize\textit{Prompt:} an old wizard with a long white beard} \\
\addlinespace[2pt]
\includegraphics[width=0.155\textwidth]{figs/kl_comparison/base/human_08_natural.png} &
\includegraphics[width=0.155\textwidth]{figs/kl_comparison/kl0/human_08_natural.png} &
\includegraphics[width=0.155\textwidth]{figs/kl_comparison/kl025/human_08_natural.png} &
\includegraphics[width=0.155\textwidth]{figs/kl_comparison/kl05/human_08_natural.png} &
\includegraphics[width=0.155\textwidth]{figs/kl_comparison/kl1/human_08_natural.png} &
\includegraphics[width=0.155\textwidth]{figs/kl_comparison/kl15/human_08_natural.png} \\
\multicolumn{6}{l}{\footnotesize\textit{Prompt:} a teenage girl with pink hair} \\
\end{tabular}
\caption{Effect of the KL-to-base weight $\lambda$ on visual style acquisition
(rank $32$). Style transfer is strong for $\lambda\!\in\![0,0.5]$ and visibly
collapses back toward the base for $\lambda\!\geq\!1$, where the adapter is
over-regularised.}
\label{fig:kl_visual}
\end{figure*}


\section{Image generation under the text adapter}
\label{app:t2i-text}
\begin{figure*}[t]
\centering
\setlength{\tabcolsep}{1.5pt}
\renewcommand{\arraystretch}{0.9}
\small
\begin{tabular}{cccccc}
\textbf{Base} & $\lambda{=}0$ & $\lambda{=}0.25$ & $\lambda{=}0.5$ & $\lambda{=}1$ & $\lambda{=}1.5$ \\
\midrule
\includegraphics[width=0.15\textwidth]{figs/t2i_text/kl0/prompt0_base.png} &
\includegraphics[width=0.15\textwidth]{figs/t2i_text/kl0/prompt0_adapter.png} &
\includegraphics[width=0.15\textwidth]{figs/t2i_text/kl025/prompt0_adapter.png} &
\includegraphics[width=0.15\textwidth]{figs/t2i_text/kl05/prompt0_adapter.png} &
\includegraphics[width=0.15\textwidth]{figs/t2i_text/kl1/prompt0_adapter.png} &
\includegraphics[width=0.15\textwidth]{figs/t2i_text/kl15/prompt0_adapter.png} \\
\multicolumn{6}{l}{\footnotesize\textit{Prompt:} a cat sitting on a chair} \\
\addlinespace[2pt]
\includegraphics[width=0.15\textwidth]{figs/t2i_text/kl0/prompt1_base.png} &
\includegraphics[width=0.15\textwidth]{figs/t2i_text/kl0/prompt1_adapter.png} &
\includegraphics[width=0.15\textwidth]{figs/t2i_text/kl025/prompt1_adapter.png} &
\includegraphics[width=0.15\textwidth]{figs/t2i_text/kl05/prompt1_adapter.png} &
\includegraphics[width=0.15\textwidth]{figs/t2i_text/kl1/prompt1_adapter.png} &
\includegraphics[width=0.15\textwidth]{figs/t2i_text/kl15/prompt1_adapter.png} \\
\multicolumn{6}{l}{\footnotesize\textit{Prompt:} a medieval castle on a hill} \\
\addlinespace[2pt]
\includegraphics[width=0.15\textwidth]{figs/t2i_text/kl0/prompt2_base.png} &
\includegraphics[width=0.15\textwidth]{figs/t2i_text/kl0/prompt2_adapter.png} &
\includegraphics[width=0.15\textwidth]{figs/t2i_text/kl025/prompt2_adapter.png} &
\includegraphics[width=0.15\textwidth]{figs/t2i_text/kl05/prompt2_adapter.png} &
\includegraphics[width=0.15\textwidth]{figs/t2i_text/kl1/prompt2_adapter.png} &
\includegraphics[width=0.15\textwidth]{figs/t2i_text/kl15/prompt2_adapter.png} \\
\multicolumn{6}{l}{\footnotesize\textit{Prompt:} a bowl of fruit on a table} \\
\addlinespace[2pt]
\includegraphics[width=0.15\textwidth]{figs/t2i_text/kl0/prompt3_base.png} &
\includegraphics[width=0.15\textwidth]{figs/t2i_text/kl0/prompt3_adapter.png} &
\includegraphics[width=0.15\textwidth]{figs/t2i_text/kl025/prompt3_adapter.png} &
\includegraphics[width=0.15\textwidth]{figs/t2i_text/kl05/prompt3_adapter.png} &
\includegraphics[width=0.15\textwidth]{figs/t2i_text/kl1/prompt3_adapter.png} &
\includegraphics[width=0.15\textwidth]{figs/t2i_text/kl15/prompt3_adapter.png} \\
\multicolumn{6}{l}{\footnotesize\textit{Prompt:} an astronaut riding a horse} \\
\end{tabular}
\caption{Image generation under the text (poetry) adapter across KL weights
(rank 16). Column 1 is the base model; columns 2 to 6 are the poetry adapter at
increasing $\lambda$, on the same prompts and seeds. The generated images stay
close to the base across the whole range.}
\label{fig:t2i_text_kl}
\end{figure*}


\section{Intervention strength and its weight-space localization}
\label{app:weight-delta}
We report the relative weight change $\lVert\Delta W\rVert_F/\lVert W\rVert_F$ that
each adapter induces on the shared backbone, aggregated over all seven projections
per layer (Tables~\ref{tab:dw-rank},~\ref{tab:dw-kl}) and resolved per module
(Fig.~\ref{fig:dw-heatmap}). The base norm $\lVert W\rVert_F$ is read from the
frozen weights; $\lVert\Delta W\rVert_F=\lVert(\alpha/r)BA\rVert_F$ is computed
from the LoRA factors via the $r\times r$ trace identity, so the full $\Delta W$
is never formed.

\begin{table}[t]\centering\small
\caption{Intervention strength per adapter: relative weight change
$\lVert\Delta W\rVert_F/\lVert W\rVert_F$ (\%, summed over all backbone
projections). It grows smoothly with rank in both directions; at matched rank
the image-generation adapter perturbs the shared backbone more than text.}
\label{tab:dw-rank}
\begin{tabular}{lcccccc}
\toprule
Rank $r$ & 2 & 4 & 8 & 16 & 32 & 64 \\
\midrule
Image (gen.)  & 0.58 & 0.75 & 0.97 & 1.25 & 1.57 & 1.91 \\
Text (poetry) & --   & 0.67 & 0.84 & 1.05 & 1.32 & 1.66 \\
\bottomrule
\end{tabular}
\end{table}

\begin{table}[t]\centering\small
\caption{KL-to-base regularizer vs.\ intervention strength
$\lVert\Delta W\rVert_F/\lVert W\rVert_F$ (\%). Larger $\lambda$ shrinks the
update, more sharply for text.}
\label{tab:dw-kl}
\begin{tabular}{lccccc}
\toprule
$\lambda$ & 0 & 0.25 & 0.5 & 1 & 1.5 \\
\midrule
Image (r32) & 1.55 & 1.58 & 1.57 & 1.45 & 1.33 \\
Text (r16)  & 1.15 & 1.15 & 1.05 & 0.89 & 0.80 \\
\bottomrule
\end{tabular}
\end{table}

\begin{figure*}[t]
\centering
\begin{subfigure}{0.32\linewidth}\centering
  \includegraphics[width=\linewidth]{figs/weight_delta/dw_heatmap_img_r16.png}
  \caption{Image gen., $r{=}16$}
\end{subfigure}\hfill
\begin{subfigure}{0.32\linewidth}\centering
  \includegraphics[width=\linewidth]{figs/weight_delta/dw_heatmap_text_r16.png}
  \caption{Text/poetry, $r{=}16$}
\end{subfigure}\hfill
\begin{subfigure}{0.32\linewidth}\centering
  \includegraphics[width=\linewidth]{figs/weight_delta/dw_heatmap_text_r64.png}
  \caption{Text/poetry, $r{=}64$}
\end{subfigure}
\caption{Per-module relative weight change
$\lVert\Delta W\rVert_F/\lVert W\rVert_F$ across layers (rows) and the seven
projections (columns). Common to both directions, the update spares the
attention query/key/value projections almost everywhere, awakening only in the
last few layers, and concentrates in the MLP. The directions differ in the fine
locus: text spreads the change broadly over \texttt{gate} across all depths,
whereas image concentrates it toward the output, peaking at the final layer's
\texttt{down} projection. Raising the rank (text $r{=}16$ vs.\ $r{=}64$) scales
the same pattern up rather than relocating it.}
\label{fig:dw-heatmap}
\end{figure*}

\section{Text-direction qualitative examples}
\label{app:text-examples}
We show base-vs-adapter continuations for the poetry text-LoRA ($r{=}16$,
$\lambda{=}0.5$); the base is the identical model with the adapter disabled, so the
only change between columns is the LoRA. Decoding is greedy with a repetition
penalty of $1.3$ (Appendix~\ref{app:impl}). Each incipit is continued from a fixed
two-line verse couplet given as hidden context to \emph{both} models (so the
unified model stays in a poem rather than its assistant persona); we display only
the text generated after the incipit. With the same primed context the adapter
sustains metrical, lineated verse, whereas the base drifts into prose or stops
short. On bare, non-poetic prompts the register does not surface --- the LoRA is a
next-token language model, not instruction-tuned.

\begin{table}[t]\centering\small
\renewcommand{\arraystretch}{1.3}
\begin{tabular}{@{}p{0.45\linewidth}p{0.45\linewidth}@{}}
\toprule
\multicolumn{2}{@{}l}{\textit{Incipit:} ``The autumn wind\,\dots''}\\
\midrule
\textbf{Base (adapter off)} & \textbf{Poetry LoRA}\\[2pt]
\footnotesize is blowing in our faces; it's a time to reflect on what has been
accomplished this year\,\dots
&
\footnotesize is blowing in our faces; it's a-blowing! --- O what shall I do?\\
\midrule
\multicolumn{2}{@{}l}{\textit{Incipit:} ``Upon the silent hill there\,\dots''}\\
\midrule
\textbf{Base (adapter off)} & \textbf{Poetry LoRA}\\[2pt]
\footnotesize lies a light of gold; --- O! how it gleams in twilight's shade: it
is not day nor night but something more divine.
&
\footnotesize stands a tree;\newline
The leaves are green as emeralds in their sleep:\newline
And on its branches hang an empty cage---\newline
A prison for some bird of sorrow's keep!\\
\midrule
\multicolumn{2}{@{}l}{\textit{Incipit:} ``Beneath the ancient oak,\,\dots''}\\
\midrule
\textbf{Base (adapter off)} & \textbf{Poetry LoRA}\\[2pt]
\footnotesize where time has worn its way; --- O how I love thee!
&
\footnotesize where I have stood;\newline
The sun was setting on a summer's day: ---\newline
I saw her face for one brief moment flash!\newline
And then she vanished from before mine eyes.\\
\midrule
\multicolumn{2}{@{}l}{\textit{Incipit:} ``O gentle heart, what\,\dots''}\\
\midrule
\textbf{Base (adapter off)} & \textbf{Poetry LoRA}\\[2pt]
\footnotesize words of love can say? --- O! for a thousand tongues to tell thee how
I feel this day; but none are fit nor true \dots\ (1 Samuel 25).
&
\footnotesize is thy secret? What art thou made of but a little child's delight in
play; \dots\ O say not ``I am nothing,'' for I see how much you have to give: your
eyes shine like stars on earth when they look at me\,\dots\\
\bottomrule
\end{tabular}
\caption{Text-direction qualitative examples (poetry LoRA, $r{=}16$,
$\lambda{=}0.5$; greedy decoding, repetition penalty $1.3$). Each incipit is
continued from a fixed verse couplet given as hidden context to \emph{both}
models; we show only the text after the incipit. With the same primed context the
adapter sustains lineated, imagistic verse (e.g.\ the emerald-leaf quatrain),
whereas the base is shorter, more abstract, or drifts into prose (top) or
scriptural citation (bottom).}
\label{tab:text-examples}
\end{table}




\section{Regularizer ablation}
\label{app:reg-ablation}
We compare the KL-to-base regularizer at fixed rank 32 in the image-generation
direction against three alternatives: L2-SP \citep{xuhong2018l2sp} (an L2 penalty
on the LoRA update, weight space), logit-MSE (a squared-error match to the base
logits, output space), and the reconstruction-loss objective of Anole
\citep{chern2024anole}, the standard recipe for fine-tuning native generators
(\texttt{CE}$+0.5\,$recon, no base anchor). All share the same training budget and
data and are evaluated identically on $80$ held-out prompts under the canonical,
trigger-free protocol used in the rest of the paper; we also report each adapter's
intervention strength $\lVert\Delta W\rVert_F/\lVert W\rVert_F$. The KL attains the
best distributional acquisition of all methods: at essentially the same
intervention strength as the unregularized baseline ($1.57$ vs.\ $1.55\%$) it
reaches a markedly lower CMMD ($.108$ vs.\ $.155$), i.e.\ it \emph{redirects} the
update rather than merely shrinking it. Logit-MSE over-constrains the output
distribution: it keeps image quality high but barely transfers the style without a
prompt trigger (DirSim $.029$, the worst of all methods), which the prompt-triggered
evaluation had masked. Weight-space L2-SP simply shrinks the update (to $0.60\%$).
The Anole reconstruction loss acquires the style comparably to the unregularized
baseline (DirSim $.162$) but with worse distributional fidelity than the KL (CMMD
$.148$ vs.\ $.108$) and, being an acquisition term rather than a base anchor, offers
no forgetting protection. The KL alone attains the best acquisition while keeping
the adapter anchored to the base. Single-seed run.

\begin{table}[t]\centering\small
\setlength{\tabcolsep}{3.5pt}
\begin{tabular}{lcccccc}
\toprule
Regularizer & $\lVert\Delta W\rVert/\lVert W\rVert$ & CMMD $\downarrow$ & KID $\downarrow$ & DINO $\uparrow$ & CLIP-IQA $\uparrow$ & DirSim $\uparrow$ \\
\midrule
None ($\lambda{=}0$) & $1.55\%$ & $.155_{\pm.009}$ & $.027_{\pm.004}$ & $.626_{\pm.013}$ & $.666_{\pm.011}$ & $.157_{\pm.007}$ \\
L2-SP                & $0.60\%$ & $.139_{\pm.009}$ & $.018_{\pm.004}$ & $.624_{\pm.011}$ & $.683_{\pm.008}$ & $\mathbf{.164}_{\pm.007}$ \\
Logit-MSE            & $1.14\%$ & $.216_{\pm.012}$ & $.106_{\pm.011}$ & $.410_{\pm.018}$ & $\mathbf{.725}_{\pm.014}$ & $.029_{\pm.006}$ \\
KL-to-base (ours)    & $1.57\%$ & $\mathbf{.108}_{\pm.010}$ & $\mathbf{.014}_{\pm.004}$ & $\mathbf{.643}_{\pm.013}$ & $.691_{\pm.010}$ & $.153_{\pm.007}$ \\
Anole (recon)       & $1.56\%$ & $.148_{\pm.009}$ & $.026_{\pm.005}$ & $.637_{\pm.011}$ & $.666_{\pm.010}$ & $.162_{\pm.007}$ \\
\bottomrule
\end{tabular}
\caption{Regularizer ablation at rank 32 (image direction, $N{=}80$, single seed;
$\pm$ one standard error; bold = best per metric). Under the canonical,
trigger-free evaluation, the KL attains the best distributional acquisition
(CMMD, KID, DINO), separating clearly from the no-regularizer baseline at matched
intervention strength ($\lVert\Delta W\rVert/\lVert W\rVert$ $1.57$ vs.\ $1.55\%$,
CMMD $.108$ vs.\ $.155$). Logit-MSE over-constrains the output distribution and
barely transfers the style without the prompt trigger (DirSim $.029$) despite high
image quality; L2-SP merely shrinks the update. The Anole reconstruction loss
acquires the style comparably to the unregularized baseline but with worse
distributional fidelity than the KL and no base anchoring.}
\label{tab:reg-ablation}
\end{table}

\clearpage
\section{Metric debugging: is the $\lambda$ optimum a bug, memorization, or noise?}
\label{app:metric-debug}
The style metrics put the image-generation acquisition optimum at $\lambda{=}0.25$,
while the eye mildly prefers $\lambda{=}0.5$. We resolve this directly by
re-scoring the \emph{already-generated} KL-sweep images (no retraining) against
both held-out and training references, and with two additional embedding metrics
(Table~\ref{tab:metric-debug}). \textbf{Not noise:} the $\lambda$-ranking of CMMD
is identical against held-out and training reals (both order
$0.25\!<\!0.5\!<\!0\!<\!1\!<\!1.5$), so the optimum is stable, not a small-$N$
artefact. \textbf{Not memorization or leakage:} for every $\lambda$ the
nearest-neighbour DINOv2 similarity to \emph{training} images is \emph{lower} than
to held-out images (and CMMD to training is higher), so the generations are not
near-copies of the training set. \textbf{A metric-choice effect, not a bug:} CMMD
(a CLIP-space distribution distance) favours $\lambda{=}0.25$, but the
memorization-robust DINOv2-space distribution distance (DINO-MMD) instead favours
$\lambda{=}0.5$ ($.035$ vs.\ $.041$), and content preservation -- the DINOv2
similarity between the adapter's image and the base image for the same prompt --
rises from $.32$ at $\lambda{=}0.25$ to $.39$ at $\lambda{=}0.5$. The CLIP style
metrics reward pure style proximity (optimum $\lambda{=}0.25$); the DINO
distribution distance and content preservation reward style transfer that keeps
the prompt content (optimum $\lambda{=}0.5$), which is what the eye weighs, and what an independent VLM judge confirms (Table~\ref{tab:judge}). This justifies the $\lambda{=}0.5$ Final recipe.

\begin{table}[h]
\centering\small
\setlength{\tabcolsep}{5pt}
\begin{tabular}{lcccccc}
\toprule
$\lambda$ & CMMD$_{\text{test}}\downarrow$ & CMMD$_{\text{train}}\downarrow$ & DINO-MMD$\downarrow$ & NN$_{\text{test}}\uparrow$ & NN$_{\text{train}}\uparrow$ & CtntPres$\uparrow$ \\
\midrule
$0$    & $.155$ & $.163$ & $.044$ & $.626$ & $.609$ & $.303$ \\
$0.25$ & $\mathbf{.100}$ & $.111$ & $.041$ & $.647$ & $.634$ & $.317$ \\
$0.5$  & $.108$ & $.124$ & $\mathbf{.035}$ & $.643$ & $.633$ & $.389$ \\
$1$    & $.205$ & $.221$ & $.091$ & $.439$ & $.411$ & $.721$ \\
$1.5$  & $.221$ & $.235$ & $.114$ & $.388$ & $.364$ & $.792$ \\
\bottomrule
\end{tabular}
\caption{Metric debugging on the KL sweep (rank 32), re-scoring the
already-generated images (no retraining). CMMD$_{\text{test/train}}$: CLIP-space
distribution distance to held-out vs.\ training reals (identical ranking $\Rightarrow$
signal not noise; train $>$ test $\Rightarrow$ no training over-fit). NN: nearest
DINOv2 similarity to held-out vs.\ training reals (train $<$ test everywhere
$\Rightarrow$ no memorization). DINO-MMD: memorization-robust DINOv2-space
distribution distance -- optimum $\lambda{=}0.5$. CtntPres: DINOv2 similarity
between the adapter and base image for the same prompt -- rises with $\lambda$. The
CLIP metrics favour $\lambda{=}0.25$; DINO-MMD and content preservation favour
$\lambda{=}0.5$.}
\label{tab:metric-debug}
\end{table}

\clearpage
\section{Medicine: qualitative before/after}
\label{app:medicine}
Same prompts, same model, adapter off (base) vs.\ on (medicine LoRA). The text
continuations show the register shift toward clinical/abstract prose; the image
grid shows that text specialisation leaves the generation direction essentially
unchanged.

\begin{table}[t]
\centering\footnotesize
\setlength{\tabcolsep}{6pt}\renewcommand{\arraystretch}{1.25}
\begin{tabular}{p{0.45\textwidth} p{0.45\textwidth}}
\toprule
\textbf{Base} & \textbf{Medicine LoRA} \\
\midrule
\multicolumn{2}{@{}p{0.95\textwidth}@{}}{\textit{Prompt:} ``What is the mechanism of action of metformin?''}\\
Metformin works by inhibiting glucose production in cells, particularly those
that are most active at high concentrations. Hello! I'm here to help you with
your question [\dots]
&
Metformin works by inhibiting glucose uptake in muscle cells and reducing
hepatic gluconeogenesis. [\dots] \\
\midrule
\multicolumn{2}{@{}p{0.95\textwidth}@{}}{\textit{Prompt:} ``A 58-year-old man with hypertension presents with acute chest pain radiating to the left arm; the most likely diagnosis is\dots''}\\
\textit{(returns another question instead of a diagnosis:)} ``What are some
common causes of sudden cardiac death in older adults?''
&
unstable angina. A coronary catheterization should be performed as soon as
possible and an appropriate treatment plan must begin immediately; if there are
no signs of myocardial infarction (MI), an electrocardiogram (ECG) and stress
test [\dots] \\
\midrule
\multicolumn{2}{@{}p{0.95\textwidth}@{}}{\textit{Prompt (general-knowledge control):} ``Who wrote the play Romeo and Juliet?''}\\
William Shakespeare. [\dots] the playwright of this famous tragedy [\dots]
($1564$--$1623$).
&
William Shakespeare. [\dots] a famous English playwright who lived between
approximately $1564$--$1616$ [\dots] including Romeo \& Juliet! \\
\bottomrule
\end{tabular}
\caption{Text continuations (greedy, deterministic; lightly trimmed, [\dots]).
On medical prompts the base stays vague or fails to answer, while the medicine
adapter gives precise clinical content and the register of PubMed (drug
mechanism, ``reducing hepatic gluconeogenesis''; a diagnosis with management). On
a general-knowledge control both answer correctly, i.e.\ the domain adaptation
does not cost general factual recall. Residual chat artifacts (``Hello!\dots'')
appear because the base is partly instruction-tuned.}
\label{tab:med-text}
\end{table}

\begin{figure}[t]
\centering
\setlength{\tabcolsep}{1.5pt}\renewcommand{\arraystretch}{1.0}
\begin{tabular}{@{}cccc@{}}
\rotatebox{90}{\footnotesize\ \ Base} &
\includegraphics[width=0.30\textwidth]{figs/med_comparison/prompt0_base.png} &
\includegraphics[width=0.30\textwidth]{figs/med_comparison/prompt1_base.png} &
\includegraphics[width=0.30\textwidth]{figs/med_comparison/prompt6_base.png} \\
\rotatebox{90}{\footnotesize Med LoRA} &
\includegraphics[width=0.30\textwidth]{figs/med_comparison/prompt0_adapter.png} &
\includegraphics[width=0.30\textwidth]{figs/med_comparison/prompt1_adapter.png} &
\includegraphics[width=0.30\textwidth]{figs/med_comparison/prompt6_adapter.png} \\
\end{tabular}
\caption{Image generation under the text--medicine LoRA: same prompts and seeds,
adapter off (top) vs.\ on (bottom). Generations are essentially unchanged,
confirming the text adapter does not degrade the generation direction (cf.\ the
near-zero Gen-CMMD).}
\label{fig:med-t2i}
\end{figure}

\section{Welsh text-direction qualitative examples}
\label{app:welsh-examples}
Base-vs-adapter continuations for the Welsh text-LoRA ($r{=}32$, $\lambda{=}0.1$);
the base is the identical model with the adapter disabled. Greedy decoding,
repetition penalty $1.3$. Two readings: (i) on \emph{Welsh} seeds the adapter
sustains Welsh while the base breaks into English/assistant text (acquisition);
(ii) on the \emph{English} seed the adapter stays in English (no bleed). As with
poetry, greedy decoding is locally repetitive; sampled decoding ($T{=}0.9$, last
row) is more fluent.

\begin{table}[t]\centering\small
\renewcommand{\arraystretch}{1.3}
\begin{tabular}{@{}p{0.45\linewidth}p{0.45\linewidth}@{}}
\toprule
\multicolumn{2}{@{}l}{\textit{Welsh seed:} ``Mae Cymru yn wlad\,\dots'' (\textit{Wales is a country})}\\
\midrule
\textbf{Base (adapter off)} & \textbf{Welsh LoRA}\\[2pt]
\footnotesize fod â'r eich, a dda i gael \dots\ Hello Mae! I am happy to help you
with your Welsh language question.
&
\footnotesize fwyaf o'r gogledd-ddinasol a ddau 10,578.234 iawn (a chyfansoddwyd
gan yr Aelodau Canolfan Llyfrau Genedlaethol). Mae hynny wedi ei bod arall am\,\dots\\
\midrule
\multicolumn{2}{@{}l}{\textit{English seed:} ``The history of Wales\,\dots''}\\
\midrule
\textbf{Base (adapter off)} & \textbf{Welsh LoRA}\\[2pt]
\footnotesize is a long and complex one, with many different cultures influencing
its development over the centuries. Hello!
&
\footnotesize is a long and complex one, with many different cultures influencing
the country over time. I will be writing about some key events in Welsh History
that have shaped our nation today\,\dots\\
\midrule
\multicolumn{2}{@{}l}{\textit{Welsh seed (sampled, $T{=}0.9$):} ``Yr oedd unwaith\,\dots'' (\textit{There was once})}\\
\midrule
\textbf{Base (adapter off)} & \textbf{Welsh LoRA}\\[2pt]
\footnotesize Hello world \dots\ Wyt yn agor fod dinio? Nochthawd, bye!
&
\footnotesize a byddai'n anhysbys, neu llawn barwas yng ngogleddol rhagolig sydd
wedi bod yn enwu am fod fel cysgu gerfansau pencampell\,\dots\\
\bottomrule
\end{tabular}
\caption{Welsh text-direction qualitative examples ($r{=}32$, $\lambda{=}0.1$). On
Welsh seeds the adapter sustains Welsh where the base slips into English/assistant
text; on the English seed it stays in English (no bleed). Greedy rows are locally
repetitive; the sampled row is more fluent.}
\label{tab:welsh-examples}
\end{table}

\begin{figure}[t]\centering
\setlength{\tabcolsep}{1.5pt}\renewcommand{\arraystretch}{0.8}\small
\begin{tabular}{cc}
\textbf{Base} & \textbf{Welsh LoRA}\\
\includegraphics[width=0.30\linewidth]{figs/t2i_text/welsh/prompt0_base.png} &
\includegraphics[width=0.30\linewidth]{figs/t2i_text/welsh/prompt0_adapter.png}\\
\multicolumn{2}{c}{\scriptsize\textit{a red fox sitting in a snowy forest at dawn}}\\
\addlinespace[1pt]
\includegraphics[width=0.30\linewidth]{figs/t2i_text/welsh/prompt2_base.png} &
\includegraphics[width=0.30\linewidth]{figs/t2i_text/welsh/prompt2_adapter.png}\\
\multicolumn{2}{c}{\scriptsize\textit{a close-up portrait of an elderly fisherman with a wool cap}}\\
\addlinespace[1pt]
\includegraphics[width=0.30\linewidth]{figs/t2i_text/welsh/prompt4_base.png} &
\includegraphics[width=0.30\linewidth]{figs/t2i_text/welsh/prompt4_adapter.png}\\
\multicolumn{2}{c}{\scriptsize\textit{a futuristic city skyline at night with neon lights}}\\
\addlinespace[1pt]
\includegraphics[width=0.30\linewidth]{figs/t2i_text/welsh/prompt6_base.png} &
\includegraphics[width=0.30\linewidth]{figs/t2i_text/welsh/prompt6_adapter.png}\\
\multicolumn{2}{c}{\scriptsize\textit{a plate of spaghetti with tomato sauce and fresh basil}}\\
\addlinespace[1pt]
\includegraphics[width=0.30\linewidth]{figs/t2i_text/welsh/prompt7_base.png} &
\includegraphics[width=0.30\linewidth]{figs/t2i_text/welsh/prompt7_adapter.png}\\
\multicolumn{2}{c}{\scriptsize\textit{an astronaut riding a horse on the surface of Mars}}\\
\end{tabular}
\caption{Image generation under the Welsh text adapter ($r{=}32$, $\lambda{=}0.1$),
base vs adapter on identical prompts and seeds. Generations stay close to the base,
consistent with the near-zero Gen-CMMD ($1.7\times10^{-3}$, Table~\ref{tab:welsh}).}
\label{fig:t2i_welsh}
\end{figure}

\end{document}
